\newpage
\begin{center}
  \textbf{\large ЗАКЛЮЧЕНИЕ}
\end{center}
\refstepcounter{chapter}
\addcontentsline{toc}{chapter}{ЗАКЛЮЧЕНИЕ}
В ходе нашего исследования была успешно разработана математическая модель для изучения динамики лазерно-индуцированных кавитационных пузырей. 
Использование метода конечных объемов позволило с высокой точностью смоделировать процессы роста и коллапса пузырей. 
Сравнение результатов моделирования с экспериментальными данными подтвердило достоверность модели, особенно в случае одиночных импульсов.

Полученные результаты способствуют углубленному пониманию явления кавитации и его применений. 
Учитывая широкий спектр применения кавитации в таких областях как неинвазивная хирургия, доставка лекарств, промышленные процессы очистки и других, понимание динамики кавитационных пузырей поможет улучшить эти технологии и повысить их безопасность и эффективность. 
В дальнейшем планируется расширение модели для учета многократных импульсов и более сложных сценариев, что позволит значительно расширить область ее применения.
Исследование выполнено при поддержке гранта Российского Научного Фонда №20-72-10161 и при инфраструктурной поддержке МГТУ им. Н.Э. Баумана.